% Introduction
% Target: 1-1.5 pages

\subsection{Motivation}
% TODO: Expand
\begin{itemize}
    \item Small LLMs (1-2B parameters) enable edge deployment but lack reasoning capabilities
    \item Emotion information improves language understanding and generation
    \item Current approaches use deep learning embeddings (WavLM, Wav2Vec2)
    \item Challenge: Are learned representations optimal for this task?
\end{itemize}

\subsection{Research Question}
\textbf{Can interpretable acoustic features based on speech science outperform deep learning embeddings for emotion-augmented knowledge distillation?}

\subsection{Our Approach}
% TODO: Expand with methodology overview
\begin{enumerate}
    \item Design 46D acoustic features: prosody (13D), spectral (30D), formants (3D)
    \item Synthesize emotion-annotated speech using TTS
    \item Train student LLM with emotion embeddings via knowledge distillation
    \item Compare acoustic features vs. WavLM across multiple scales
\end{enumerate}

\subsection{Key Contributions}
\begin{enumerate}
    \item \textbf{Novel Feature Design:} 46-dimensional interpretable acoustic features
    \item \textbf{Empirical Discovery:} 40.45\% improvement (vs. 27.82\% for WavLM)
    \item \textbf{Efficiency Analysis:} 8.06$\times$ dimensional efficiency
    \item \textbf{Theoretical Insight:} Speech science > deep learning for emotion encoding
    \item \textbf{Scale Validation:} Tested on 100, 500, and 1000 samples
\end{enumerate}

\subsection{Paper Organization}
% TODO: Add after finalizing structure
Section~\ref{sec:related} reviews related work.
Section~\ref{sec:method} describes our acoustic feature design and training methodology.
Section~\ref{sec:setup} details experimental configuration.
Section~\ref{sec:results} presents main results.
Section~\ref{sec:analysis} provides feature importance and ablation studies.
Section~\ref{sec:discussion} discusses implications and limitations.
